% DOCX paragraph 1352
\chapter*{El final del viaje}\label{conclusion}
\addcontentsline{toc}{chapter}{El final del viaje}
\markboth{El final del viaje}{El final del viaje}

% DOCX paragraph 1354
La gravedad\ldots{} ¡atrae! Esto era obvio antes de que el lector empezara este libro, y lo es aún más ahora que ha llegado al final. Al mismo tiempo, sin embargo, espero que ahora esté de acuerdo conmigo en que la gravedad también es \emph{atractiva}, lo cual resulta mucho menos obvio.

% DOCX paragraph 1355
Este libro se escribió con la finalidad de mostrar cuán atractiva es la gravedad y cómo está inextricablemente ligada a lo que sabemos del mundo que nos rodea y de los rincones más remotos del universo, que también se rigen por sus leyes.

% DOCX paragraph 1356
Juntos hemos completado un viaje virtual al mundo de la gravedad, tal como nos lo reveló Einstein hace más de cien años. Puedo imaginar que ha sido un viaje agotador para algunos lectores: los he obligado a ejercitar su imaginación de manera extraordinaria y los he expuesto a conceptos que, lo reconozco, son difíciles de aceptar, al menos cuando nos los encontramos por primera vez. A mí me ocurrió lo mismo cuando, a los 20 años, empecé a estudiar esta teoría por mi cuenta. Pero ¡que nadie diga que no avisé de que hacían falta abundante imaginación y una buena dosis de paciencia!

% DOCX paragraph 1357
En este viaje hemos buscado respuestas a una serie de preguntas sencillas, aunque no triviales, que me gustaría recordar aquí:

% DOCX paragraph 1359
\emph{¿Por qué una manzana cae del árbol en vez de flotar en el espacio? ¿Qué es el espaciotiempo? ¿En qué consiste su curvatura y cómo se produce? ¿Se puede doblar el tiempo? ¿Cómo funciona un agujero negro y cómo podemos «construir» uno? ¿Cómo es posible fotografiarlo si no emite luz? ¿Qué son las ondas gravitatorias y por qué son difíciles de medir?}

% DOCX paragraph 1361
Quizá estas preguntas sigan pareciendo absurdas y sin respuesta. O tal vez no; tal vez los lectores las encuentren ahora razonables e incluso interesantes. Por último, quizá tengan para cada una una respuesta que, después de todo, tenga sentido. Y, si ese es el caso, el esfuerzo de escribir este libro, ¡y el de los lectores al leerlo!, no habrá sido en vano.

% DOCX paragraph 1362
Aquí termina el viaje que hemos hecho juntos, al menos de momento. Confío en que, al cerrar el libro, quede una huella del mensaje más importante que he intentado transmitir: no dudar en usar la curiosidad y hacerse preguntas, incluso preguntas que puedan parecer extrañas. Que la intuición sea tu guía, pero no tu maestro. Como he mostrado con innumerables ejemplos, nuestra idea de cómo «funciona» el mundo está ligada inevitablemente a la experiencia que tenemos de él, que es necesariamente limitada. Lo que experimentamos a lo largo de nuestra vida en la Tierra no es más que una gota en el océano de lo físicamente posible.

% DOCX paragraph 1363
Por último, y por encima de todo, no limites tu imaginación. Antes que las matemáticas avanzadas, las simulaciones complejas y los experimentos sofisticados, todos ellos indispensables para el progreso de la ciencia, están la agilidad de nuestra mente y los viajes de la imaginación. Más que ninguna otra cosa, eso es lo que nos permite ampliar los límites del conocimiento.

% DOCX paragraph 1366
Fráncfort, 25 de septiembre de 2020
